\begin{SCn}
\begin{small}

\scnheader{Google Cloud AI}
\scnidtf{Облачная платформа искусственного интеллекта от Google}
\scniselement{инструмент для разработки ИИ-решений}
\begin{scnindent}
\scnidtf{MLaaS (Machine Learning as a Service)}
\scnidtf{Платформа для генеративного ИИ}
\end{scnindent}
\scnrelfrom{разработчик}{Google}
\scnrelfrom{год создания}{2016}
\scntext{назначение}{Предоставление комплексных инструментов для создания, обучения, развертывания и управления ИИ-моделями, включая обработку мультимодальных данных, прогнозную аналитику и автоматизацию бизнес-процессов}

\begin{scnrelfromset}{архитектурные особенности}
\scnfileitem{Единая платформа Vertex AI}
\begin{scnindent}
\scntext{уточнение}{Объединяет конвейеры ML (Kubeflow), AutoML и MLOps в едином интерфейсе}
\end{scnindent}
\scnfileitem{AI Hypercomputer}
\begin{scnindent}
\scntext{уточнение}{Комбинирует TPU (Tensor Processing Units), GPU (NVIDIA Blackwell) и CPU для обработки data-intensive workloads}
\end{scnindent}
\scnfileitem{Интеграция с Google Ecosystem}
\begin{scnindent}
\scntext{пример}{Автоматическая синхронизация с BigQuery, Cloud Storage, AlloyDB и Kubernetes Engine}
\end{scnindent}
\scnfileitem{Поддержка гибридных развертываний}
\begin{scnindent}
\scntext{уточнение}{Интеграция с локальными системами через Anthos}
\end{scnindent}
\end{scnrelfromset}

\begin{scnrelfromset}{функциональные возможности}
\scnfileitem{возможность интеграции с генеративным ИИ (Gemini)}
\scnfileitem{возможность интеграции с Imagen}
\begin{scnindent}
\scntext{назначение}{Генерация фотореалистичных изображений по текстовым промптам}
\end{scnindent}
\scnfileitem{возможность интеграции с Veo}
\begin{scnindent}
\scntext{назначение}{Создание HD-видео в кинематографических стилях}
\end{scnindent}
\scnfileitem{возможность интеграции с Code Assist}
\begin{scnindent}
\scntext{назначение}{Генерация и оптимизация кода с контекстным анализом}
\end{scnindent}
\scnfileitem{возможность интеграции с Agentspace}
\begin{scnindent}
\scntext{назначение}{Поиск информации в корпоративных данных через чат-интерфейс}
\end{scnindent}


\scnfileitem{возможность использования мультиагентных систем}

\scnfileitem{возможность простой обработки данных}
\begin{scnindent}
\scntext{уточнение}{Speech-to-Text, Cloud Translation, Vision AI, Healthcare Data Engine}
\end{scnindent}

\scnfileitem{MLOps и кастомизация}

\scnfileitem{Model Garden}
\begin{scnindent}
\scntext{назначение}{Каталог 200+ моделей (Gemini, Stable Diffusion, BERT)}
\end{scnindent}
\scnfileitem{Vertex AI Pipelines}
\begin{scnindent}
\scntext{назначение}{Оркестрация рабочих процессов ML с мониторингом}
\end{scnindent}
\scnfileitem{MedLM}
\begin{scnindent}
\scntext{назначение}{Специализированные модели для медицины (анализ рентгенограмм)}
\end{scnindent}

\scnfileitem{возможность интеграции BigQuery}
\begin{scnindent}
\scntext{назначение}{SQL-аналитика с встроенными ML-моделями}
\end{scnindent}
\scnfileitem{возможность интеграции Workspace}
\begin{scnindent}
\scntext{назначение}{Генерация контента в Gmail, Документах, Таблицах}
\end{scnindent}
\scnfileitem{возможность интеграции Looker}
\begin{scnindent}
\scntext{назначение}{Визуализация данных с ИИ-инсайтами}
\end{scnindent}
\end{scnrelfromset}

\begin{scnrelfromset}{преимущества}
\scnfileitem{Скорость внедрения}

\scnfileitem{Стартапы}
\begin{scnindent}
\scntext{пример}{Развертывание приложений за 3 недели}
\end{scnindent}

\scnfileitem{Масштабируемость}
\begin{scnindent}
\scntext{уточнение}{Поддержка кластеров до 15,000 узлов (GKE) и 500 млн триплетов данных}
\end{scnindent}

\scnfileitem{Безопасность}
\begin{scnindent}
\scntext{уточнение}{Шифрование данных и VPC-контроль, хранение данных в 23+ регионах (GDPR/HIPAA), защита конфиденциальной информации}
\end{scnindent}

\scnfileitem{развитая экосистема}

\scnfileitem{партнерства с различными компаниями}
\end{scnrelfromset}

\begin{scnrelfromset}{недостатки}
\scnfileitem{Стоимость}
\begin{scnindent}
\scntext{уточнение}{Расширенные функции TPU/GPU требуют высоких затрат}
\end{scnindent}

\scnfileitem{Сложность миграции}
\begin{scnindent}
\scntext{пример}{Трудности переноса проектов с Azure ML}
\end{scnindent}

\scnfileitem{Ограничения кастомизации}
\begin{scnindent}
\scntext{уточнение}{Низкоуровневая настройка моделей требует экспертизы}
\end{scnindent}
\end{scnrelfromset}

\begin{scnrelfromset}{библиографические источники}
\scnfileitem{Google Cloud AI: Official Documentation // https://cloud.google.com/ai. Дата обращения: 27.05.2025}
\scnfileitem{Мультиагентные системы в Vertex AI // Habr. 2024. https://habr.com/ru/companies/bothub/news/899870. Дата обращения: 27.05.2025}
\scnfileitem{AI Tools for Healthcare // DSMedia. 2024. https://dsmedia.pro/news/google-cloud-predstavljaet-novye-instrumenty-iskusstvennogo-intellekta-dlja-zdravoohranenija. Дата обращения: 27.05.2025}
\scnfileitem{Gemini in Workspace // Google. https://workspace.google.com/intl/ru/solutions/ai. Дата обращения: 27.05.2025}
\end{scnrelfromset}

%---------------------------------------

\scnheader{Naumen KMS}
\scnidtf{Система управления знаниями российской разработки}
\scniselement{корпоративное решение для клиентского сервиса}
\begin{scnindent}
\scnidtf{платформа Knowledge Management System}
\scnidtf{импортозамещающий аналог Confluence и SharePoint}
\end{scnindent}
\scnrelfrom{разработчик}{Naumen}
\scnrelfrom{год создания}{2020-е}
\scntext{назначение}{Централизация корпоративных знаний для ускорения обработки запросов, обучения сотрудников и обеспечения соответствия регуляторным требованиям РФ}

\begin{scnrelfromset}{архитектурные особенности}
\scnfileitem{On-premise/частное облако}
\begin{scnindent}
\scntext{уточнение}{Поддержка российских ОС: «Альт СП», «РЕД ОС»}
\end{scnindent}
\scnfileitem{Модульная экосистема}
\begin{scnindent}
\scntext{интеграции}{Совместимость с Naumen Contact Center, Naumen Service Desk, ERP-системами}
\end{scnindent}
\scnfileitem{Безопасность данных}
\begin{scnindent}
\scntext{соответствие}{ФЗ-152, ГОСТ Р 57580}
\end{scnindent}
\end{scnrelfromset}

\begin{scnrelfromset}{функциональные возможности}
\scnfileitem{Генеративный ИИ (GPT AssistTool)}
\scnfileitem{Автогенерация ответов}
\begin{scnindent}
\scntext{уточнение}{Формирование ответов на основе статей БЗ со ссылками на источники}
\end{scnindent}
\scnfileitem{Контекстный поиск}
\begin{scnindent}
\scntext{назначение}{Анализ семантики запросов операторов}
\end{scnindent}

\scnfileitem{Управление контентом}
\scnfileitem{Визуальный редактор}
\begin{scnindent}
\scntext{назначение}{Поддержка шаблонов для артикулов, этикеток, инструкций}
\end{scnindent}
\scnfileitem{Версионность}
\begin{scnindent}
\scntext{уточнение}{История изменений и восстановление предыдущих редакций}
\end{scnindent}
\scnfileitem{Миграция данных}
\begin{scnindent}
\scntext{уточнение}{Автоматический перенос контента из Confluence/XWiki без потерь}
\end{scnindent}

\scnfileitem{Аналитика знаний}
\scnfileitem{Отчёты по эффективности}
\begin{scnindent}
\scntext{уточнение}{Анализ популярности статей, времени поиска, пробелов в знаниях}
\end{scnindent}
\scnfileitem{Персонализированные ленты}
\begin{scnindent}
\scntext{назначение}{Уведомления об изменениях релевантных разделов}
\end{scnindent}

\scnfileitem{Интеграция с бизнес-процессами}
\scnfileitem{Скрипты для контакт-центров}
\begin{scnindent}
\scntext{учочнение}{Готовые диалоговые сценарии}
\end{scnindent}
\scnfileitem{Тестирование компетенций}
\begin{scnindent}
\scntext{назначение}{Конструктор тестов для сотрудников с интеграцией в Naumen LMS}
\end{scnindent}
\scnfileitem{Ролевая модель доступа}
\begin{scnindent}
\scntext{назначение}{Гибкие права для отделов/внешних партнеров}
\end{scnindent}

\scnfileitem{Виджеты и автоматизация}
\scnfileitem{Виджет быстрого ответа}
\begin{scnindent}
\scntext{назначение}{Встраивание в интерфейсы за 5 минут}
\end{scnindent}
\scnfileitem{Интеллектуальные уведомления}
\begin{scnindent}
\scntext{уточнение}{Оповещения о новых регламентах или изменениях продуктов}
\end{scnindent}
\end{scnrelfromset}

\begin{scnrelfromset}{преимущества}
\scnfileitem{Импортозамещение}
\begin{scnindent}
\scntext{подтверждение}{Лидер рейтинга Call Center Guru (2021), включение в реестр отечественного ПО}
\end{scnindent}

\scnfileitem{Экономическая эффективность}
\scnfileitem{Экономия времени}
\begin{scnindent}
\scntext{пример}{Совкомбанк таким образом экономит очень много операторского времени}
\end{scnindent}

\scnfileitem{Адаптивность}
\begin{scnindent}
\scntext{уточнение}{Готовые шаблоны для ритейла, банков, телекома}
\end{scnindent}
\end{scnrelfromset}

\begin{scnrelfromset}{недостатки}
\scnfileitem{Отраслевая специализация}
\begin{scnindent}
\scntext{ограничение}{Фокус на корпоративный сектор, менее гибка для R\&D-проектов}
\end{scnindent}

\scnfileitem{Сложность миграции}
\begin{scnindent}
\scntext{пример}{Необходимость ручной доработки шаблонов при переносе из Notion}
\end{scnindent}
\end{scnrelfromset}

\begin{scnrelfromset}{библиографические источники}
\scnfileitem{Официальный сайт Naumen KMS // https://www.naumen.ru/products/kms/. Дата обращения: 27.05.2025}
\scnfileitem{Naumen KMS возглавил рейтинг Call Center Guru // ComNews. 2021. https://www.comnews.ru/content/214812/2021-06-03/2021-w22/naumen-kms-vozglavil-reyting. Дата обращения: 27.05.2025}
\scnfileitem{Кейс ритейлера: Оптимизация бизнес-показателей // Naumen Academy. https://www.naumen.ru/products/kms/academy/articles/kak-baza-znaniy. Дата обращения: 27.05.2025}
\scnfileitem{Новости о выходе на рынок // NAUMEN. https://www.naumen.ru/events/news/4743/. Дата обращения: 27.05.2025}
\end{scnrelfromset}

%=======================================

\scnheader{Сравнение Google Cloud AI и Naumen KMS}
\begin{scneqtovector}
    \scnitem{Google Cloud AI}
    \scnitem{Naumen KMS}
\end{scneqtovector}

\begin{scnrelfromset}{сходства}
    \scnfileitem{Поддержка генеративного ИИ}
    \begin{scnindent}
        \scntext{детализация}{Gemini (Google) и GPT AssistTool (Naumen) для создания контента/ответов}
    \end{scnindent}
    \scnfileitem{Интеграция с бизнес-системами}
    \begin{scnindent}
        \scntext{примеры}{BigQuery/Looker (Google) ↔ ERP/Contact Center (Naumen)}
    \end{scnindent}
    \scnfileitem{Обработка естественного языка}
    \begin{scnindent}
        \scntext{уточнение}{Мультиязычный NLP (Google) vs русскоязычная оптимизация (Naumen)}
    \end{scnindent}
    \scnfileitem{Аналитика знаний}
\end{scnrelfromset}

\begin{scnrelfromset}{различия}
    \scnfileitem{Архитектура: Google Cloud AI - публичное облако (Vertex AI + TPU/GPU), Naumen KMS - on-premise/частное облако (поддержка «Альт СП»/«РЕД ОС»)}
    \scnfileitem{ML-подход: Google Cloud AI использует предобученные модели (Gemini, MedLM) + тонкую настройку, Naumen KMS обучает GPT AssistTool на внутренних базах знаний}
    \scnfileitem{Безопасность: Google Cloud AI соответствует глобальным стандартам (GDPR, HIPAA), Naumen KMS соответствует ФЗ-152 и ГОСТ Р 57580}
    \scnfileitem{Масштабируемость: Google Cloud AI поддерживает кластеры до 15,000 узлов (GKE) и 500M+ триплетов, Naumen KMS масштабируется до 3,000 пользователей (кейс Совкомбанка)}
    \scnfileitem{Целевые сценарии: Google Cloud AI для стартапов, мультиязычных проектов и Big Data, Naumen KMS для корпораций РФ, госсектора и импортозамещения}
    \scnfileitem{Экономика: Google Cloud AI имеет высокую стоимость TPU/GPU, Naumen KMS обеспечивает экономию 17.5 млн минут/год операторского времени}
\end{scnrelfromset}

\begin{scnrelfromset}{рекомендации по выбору}
    \scnfileitem{Выбрать Google Cloud AI если требуется мультимодальный ИИ (текст+аудио+видео)}
    \scnfileitem{Выбрать Google Cloud AI если проект включает мультиязычные данные (>135 языков)}
    \scnfileitem{Выбрать Google Cloud AI если критична скорость развертывания (приложения за 3 недели)}
    \scnfileitem{Выбрать Google Cloud AI если нужна работа с Big Data через BigQuery ML}
    
    \scnfileitem{Выбрать Naumen KMS если обязательно соответствие ФЗ-152/ГОСТ Р 57580}
    \scnfileitem{Выбрать Naumen KMS если требуется импортозамещение Confluence/SharePoint}
    \scnfileitem{Выбрать Naumen KMS если фокус на автоматизацию контакт-центров (сокращение AHT на 180 сек)}
    \scnfileitem{Выбрать Naumen KMS если нужна интеграция с legacy-системами (Битрикс24, 1С)}
    
    \scnfileitem{Гибридный сценарий: Использовать Google Cloud AI для R\&D (генерация ML-моделей), а Naumen KMS - для развертывания в production с соблюдением регуляторики РФ}
\end{scnrelfromset}

\end{small}
\end{SCn}